\chapter{Diseño e implementación del modelo a desarrollar}
\label{cap:diseno}

Este capítulo articula el modelo de investigación y su materialización en software. Se presenta primero el \emph{diseño metodológico} —el tipo de investigación y el ciclo de la investigación basada en diseño que este documento cubre, las variables y su matriz de consistencia, la población, la muestra y los casos de estudio, el procedimiento con sus instrumentos numéricos y didácticos, y los criterios de aceptación con que se contrasta la hipótesis— y, a continuación, el software EduFEM, que concreta ese modelo.

\section{Diseño metodológico}
\label{sec:metodologia}

La Introducción presentó una visión general del enfoque de investigación; esta sección lo desarrolla siguiendo la estructura de un modelo de investigación: la conceptualización y el alcance del modelo, la definición de las variables, el procedimiento y la toma de observaciones, la validación de esas observaciones y el contraste de la hipótesis con los resultados.

\subsection{Tipo de investigación y ciclo de la investigación basada en diseño}
\label{sec:proceso-met}

Por su finalidad, el trabajo es de carácter \emph{propositivo} —la propuesta de un artefacto que resuelve un problema formativo—; por su naturaleza es investigación aplicada de tipo \emph{tecnológico}, la que produce conocimiento operativo para transformar una realidad concreta \autocite[pp.~90-92]{garciacordoba2005tecnologica}; por el abordaje del problema, el enfoque es \emph{mixto}: cuantitativo en la verificación y validación del motor, donde la corrección del artefacto se establece midiendo errores numéricos, tensiones, desplazamientos y tasas de convergencia frente a referencias de exactitud conocida \autocite[p.~6]{sampieri2018metodologia}, y cualitativo, por análisis documental, en la evaluación del diseño, donde cada decisión se confronta con los principios de aprendizaje y los criterios de calidad de software educativo del \autoref{cap:marco_teorico}; y por su nivel es \emph{descriptiva} y \emph{comparativa}: caracteriza el comportamiento numérico del motor, contrasta de forma sistemática el desempeño de los elementos Q4 y Q9 y describe el grado en que el artefacto encarna la exposición que la hipótesis postula.

El modelo que sostiene la investigación es el de la \emph{investigación basada en diseño}, presentado en la \autoref{sec:evaluacion-software-educativo}: el estudio sistemático del diseño, el desarrollo y la evaluación de una intervención educativa como solución a un problema de la práctica, que avanza por ciclos de prototipo y evaluación y en el que los criterios de calidad de la intervención —relevancia, consistencia, practicidad y efectividad— se evalúan de forma progresiva \autocite{plomp2013edr,nieveen2013formative}. Este documento reporta el \emph{primer ciclo}: la construcción del artefacto y la evaluación de su relevancia, su consistencia y su practicidad esperada, más la corrección del contenido que expone; la efectividad, que exige medir la comprensión en estudiantes, corresponde al segundo ciclo y se deja diseñada en las recomendaciones. En consecuencia, la unidad de análisis de este ciclo es el artefacto y no el estudiante: las variables del problema científico se observan en él con los instrumentos de la \autoref{sec:validacion-datos}. El objeto de estudio declarado en la Introducción —el proceso de enseñanza y aprendizaje del análisis de medios continuos en elasticidad plana por el MEF— se estudia a través del medio didáctico que el campo de acción propone, y ese medio se contrasta, en lo pedagógico, con la literatura y, en lo numérico, con un \emph{modelo de simulación numérica}: una representación ideal del objeto en la que el investigador abstrae y sistematiza las relaciones que considera esenciales \autocite[p.~21]{alvarez_metodologia}, ejercitada aquí por simulación computacional, resolviendo problemas del dominio con el motor desarrollado y contrastando sus resultados con referencias de exactitud conocida. En ese experimento numérico no hay unidades experimentales, tratamientos ni aleatorización: el motor es determinista, de modo que la confiabilidad no descansa en réplicas sino en la reproducibilidad exacta de cada corrida y en la convergencia de sus resultados bajo refinamiento de malla; la incertidumbre numérica se cuantifica con las tasas de convergencia y, cuando la referencia no es exacta, con la extrapolación de Richardson de la propia secuencia de mallas \autocite[pp.~309-312]{oberkampf2010vv}.

El trabajo combinó dos ciclos articulados de desarrollo. El \emph{desarrollo del software} siguió un proceso iterativo e incremental: cada componente del canal de cálculo del MEF (\autoref{cap:marco_teorico}) se implementó y se sometió a prueba de regresión antes de integrarse al conjunto. La \emph{evaluación} se realizó en dos frentes: la corrección del motor, mediante verificación de código por el método de soluciones manufacturadas y contrastación con dos referencias externas, la solución analítica de la viga de Timoshenko y la membrana de Cook \autocite{oberkampf2010vv,salari2000mms}; y el potencial didáctico del artefacto, mediante los instrumentos documentales de la \autoref{sec:validacion-datos}, construidos sobre la versión del software que acompaña a este documento.

\subsection{Variables e indicadores}
\label{sec:variables-met}

La \autoref{tab:variables} operacionaliza las variables del estudio, es decir, especifica para cada una las operaciones concretas con que se la mide \autocite[p.~137]{sampieri2018metodologia}. La \emph{variable independiente} es la forma en que la herramienta expone el canal de cálculo, con tres dimensiones: transparencia, interactividad y verificabilidad numérica. Se observa en el artefacto con indicadores contables: etapas del canal con módulo interactivo y con desarrollo numérico en la memoria, vínculo entre el lienzo y el módulo y conmutador entre fórmula y valor, y coincidencia entre lo que el software expone y lo que su motor verificado calcula. La \emph{variable dependiente} es la comprensión de los fundamentos del MEF por el estudiante. De ella, este ciclo observa la dimensión que puede establecerse sobre el artefacto, su \emph{potencial didáctico} —el grado en que, según el diseño y la literatura, el artefacto está en condiciones de producir el efecto buscado—, con tres subdimensiones tomadas de los criterios de calidad de una intervención educativa \autocite{nieveen2013formative}: la relevancia, medida por la cobertura de las unidades de contenido y los niveles cognitivos en la tabla de especificaciones; la consistencia, medida por la proporción de principios de aprendizaje con encarnación verificable en la matriz de trazabilidad; y la practicidad esperada, medida por las dimensiones del instrumento de evaluación heurística con evidencia. La dimensión restante, el \emph{aprendizaje logrado}, se mide por la ganancia entre una preprueba y una posprueba de conocimientos, corresponde al segundo ciclo y figura en la tabla como no medida en este trabajo. La \emph{variable interviniente} es EduFEM, la solución tentativa con que se manipula la independiente. Las magnitudes del experimento numérico con que se establece la corrección del motor no son variables del problema científico, y la tabla las distingue con ese nombre: los \emph{factores} —la definición del modelo estructural (geometría, material, acciones y apoyos) y las decisiones de discretización, el tipo de elemento (Q4 o Q9) y la densidad de malla (tamaño característico $h$, o $N\times N$)— y las \emph{magnitudes medidas} —la respuesta estructural calculada (tensiones y desplazamientos) junto con su exactitud frente a las soluciones de referencia, medida por las normas de error $L^2$ y $H^1$ del desplazamiento, por la norma $L^2$ del campo de tensiones recuperado y por los errores relativos en tensión normal y en flecha—. El tipo de análisis (tensión o deformación plana) y el orden de cuadratura se mantienen como \emph{controlados}, y el desempeño del solucionador se registra como \emph{observado}, sin criterio de aceptación, para documentar la escalabilidad del motor.

\begin{table}[htbp]
\centering
\caption{Operacionalización de las variables de la investigación.}
\label{tab:variables}
{\footnotesize
\setlength{\tabcolsep}{3pt}
\renewcommand{\arraystretch}{1.15}
\begin{tabular}{@{}>{\raggedright\arraybackslash}p{2.9cm}>{\raggedright\arraybackslash}p{3.0cm}>{\raggedright\arraybackslash}p{5.0cm}>{\raggedright\arraybackslash}p{2.7cm}@{}}
\toprule
\textbf{Variable} & \textbf{Dimensión} & \textbf{Indicador} & \textbf{Instrumento} \\
\midrule
Independiente: exposición del canal de cálculo & Transparencia; interactividad; verificabilidad numérica & Etapas con módulo interactivo y con desarrollo numérico en la memoria; vínculo lienzo--módulo y conmutador fórmula--valor; coincidencia entre lo expuesto y lo calculado por el motor verificado & Inventario de módulos; memoria del ejemplo canónico; pruebas de generación de la memoria \\
Dependiente: comprensión de los fundamentos del MEF & Potencial didáctico: relevancia & Unidades de contenido cubiertas por nivel cognitivo & Tabla de especificaciones \\
 & Potencial didáctico: consistencia & Principios de aprendizaje con encarnación verificable & Mapa de conjeturas y matriz de trazabilidad \\
 & Potencial didáctico: practicidad esperada & Dimensiones del instrumento con evidencia & Evaluación heurística (LORI) \\
 & Aprendizaje logrado (no medido en este ciclo) & Ganancia entre preprueba y posprueba & Prueba de conocimientos (segundo ciclo) \\
Interviniente: EduFEM & Versión evaluada & Software distribuido con este documento & Instalador y repositorio \\
\midrule
Factores del experimento numérico & Definición del modelo; discretización & Dominio y malla; $E$, $\nu$, espesor $t$; cargas y restricciones; Q4/Q9; tamaño característico $h$ (o $N\times N$) y número de GDL & Guion de cada caso \\
Magnitudes medidas & Respuesta estructural; exactitud & $\sigma_x,\sigma_y,\tau_{xy},\sigma_{\mathrm{VM}}$; $u,v$ nodales; tasa de convergencia (orden $p$); normas $L^2$ y $H^1$ del desplazamiento y $L^2$ de $\sigma^*$; error relativo (\%) en $\sigma_x$ y en flecha frente al analítico y a SAP2000; residuo de la suma de reacciones & Solucionador del MEF; guiones de V\&V \\
Controlados & Tipo de análisis; integración & Tensión o deformación plana; orden de cuadratura & Constantes del guion de cada caso \\
Observado: desempeño & Tiempo; memoria & Tiempos de ensamblaje y de solución; memoria de $\bm{K}$ densa y dispersa & Guion de medición de tiempos \\
\bottomrule
\end{tabular}}
\end{table}

\subsection{Matriz de consistencia}
\label{sec:matriz-consistencia}

El \textbf{problema científico}, en la formulación única de la Introducción —¿de qué manera la exposición transparente, interactiva y numéricamente verificable del canal de cálculo completo del MEF en elasticidad plana, mediante un software educativo que opera sobre el modelo del propio estudiante, apoya la comprensión de los fundamentos del método en los estudiantes de ingeniería civil?—, se atiende mediante el \textbf{objetivo general} de desarrollar EduFEM y se contrasta con la \textbf{hipótesis} de que esa exposición apoyará la comprensión, cuya parte comprobable en este ciclo son las cuatro cláusulas de relevancia, corrección, consistencia y practicidad esperada. La \autoref{tab:consistencia} traza la correspondencia entre cada objetivo específico, la pregunta de investigación que responde, las variables o indicadores asociados, el método o instrumento empleado y la evidencia que lo respalda en el documento.

\begin{table}[htbp]
\centering
\caption[Matriz de consistencia de la investigación.]{Matriz de consistencia: correspondencia entre objetivos específicos, preguntas de investigación, variables, métodos y evidencia.}
\label{tab:consistencia}
{\footnotesize
\setlength{\tabcolsep}{3pt}
\renewcommand{\arraystretch}{1.15}
\begin{tabular}{@{}>{\raggedright\arraybackslash}p{3.0cm}>{\raggedright\arraybackslash}p{1.6cm}>{\raggedright\arraybackslash}p{3.1cm}>{\raggedright\arraybackslash}p{3.3cm}>{\raggedright\arraybackslash}p{2.6cm}@{}}
\toprule
\textbf{Objetivo específico} & \textbf{Pregunta asociada} & \textbf{Variable / indicador} & \textbf{Método / instrumento} & \textbf{Evidencia} \\
\midrule
Fundamentar teóricamente el MEF, los principios de aprendizaje y los criterios de evaluación & PI-1 a PI-4 & Canal de cálculo cubierto; principios y criterios adoptados & Revisión de la literatura clásica del MEF, de la V\&V, de la enseñanza del método y de la evaluación de software educativo & \autoref{cap:marco_teorico} \\
Implementar el motor MEF 2D Q4/Q9 & PI-2 & Tipo de elemento; exactitud (normas $L^2$, $H^1$ y $L^2$ de $\sigma^*$) & Implementación pura en NumPy/SciPy; pruebas de regresión & \S\ref{sec:motor}; \autoref{cap:resultados} (MMS) \\
Diseñar la interfaz pre/proceso/post & PI-3, PI-4 & Fases que comparten el lienzo; sincronización lienzo--tabla; dimensiones de LORI & Diseño centrado en un lienzo único compartido; evaluación heurística & \S\ref{sec:pre-proceso}--\S\ref{sec:post-proceso}; \autoref{tab:fases}; \autoref{tab:lori} \\
Desarrollar los módulos educativos & PI-1, PI-3 & Etapas del canal con módulo interactivo; principios encarnados & Capas interactivas superpuestas a la malla real; inventario de módulos; matriz de trazabilidad & \S\ref{sec:educacion}; \autoref{tab:especificaciones}; \autoref{tab:trazabilidad} \\
Verificar y validar el motor & PI-2 & Exactitud; tasa de convergencia; bloqueo por cortante & MMS, viga de Timoshenko, membrana de Cook & \autoref{cap:resultados} \\
Generar la memoria de cálculo e interoperabilidad & PI-1, PI-3 & Etapas con desarrollo numérico en la memoria; efecto del ejemplo resuelto encarnado; formatos con ida y vuelta & Generación de PDF con las funciones del motor; importación y exportación DXF y CSV & \S\ref{sec:memoria-io}; \autoref{anexo:memoria}; \autoref{sec:consistencia-interna} \\
Evaluar el potencial didáctico y diseñar la prueba de campo & PI-1, PI-3, PI-4 & Unidades cubiertas por nivel; principios con encarnación; dimensiones con evidencia & Tabla de especificaciones; mapa de conjeturas y matriz de trazabilidad; LORI & \autoref{sec:cobertura}; \autoref{sec:validacion-hipotesis}; recomendaciones \\
\bottomrule
\end{tabular}}
\end{table}

Las preguntas de investigación referidas en la matriz son las enunciadas en la Introducción, una por cláusula de la hipótesis: \textbf{PI-1}, sobre la relevancia, es decir, la cobertura de los contenidos y los niveles cognitivos; \textbf{PI-2}, sobre la corrección del motor frente a soluciones analíticas y comerciales y la reproducción de los fenómenos numéricos; \textbf{PI-3}, sobre la consistencia entre cada decisión de diseño y un principio de aprendizaje documentado; y \textbf{PI-4}, sobre la practicidad esperada según los criterios de calidad evaluables sin usuarios.

El primer objetivo específico —la fundamentación teórica— es transversal: sustenta las cuatro preguntas de investigación y se materializa en el \autoref{cap:marco_teorico}, de ahí que la matriz le asocie las cuatro. El séptimo objetivo cierra la matriz: sus tres instrumentos observan la variable dependiente en el artefacto, y la prueba de campo que deja diseñada es la que observará esa variable en el estudiante. El sexto objetivo responde a \textbf{PI-1} y \textbf{PI-3} en su parte principal: la memoria de cálculo es la forma escrita en que el canal se expone de manera transparente y contrastable, cubre las nueve etapas en la tabla de especificaciones y es la encarnación del efecto del ejemplo resuelto en la matriz de trazabilidad. La interoperabilidad de datos (DXF, CSV) es su complemento instrumental: no deriva de una pregunta de investigación, pero se verifica por ida y vuelta e idempotencia (\autoref{sec:consistencia-interna}) para que el modelo del alumno pueda entrar y salir de la herramienta sin pérdida.

\subsection{Población, muestra y casos de estudio}
\label{sec:poblacion}

La \emph{población} del problema científico son los estudiantes de la Carrera de Ingeniería Civil de la Universidad Autónoma Tomás Frías que cursan las asignaturas en que se introduce el análisis por el MEF. % DATO PENDIENTE: nombre, sigla y semestre de la asignatura segun el plan de estudios vigente
Este ciclo no toma una muestra de esa población, porque su unidad de análisis es el artefacto: la comprensión no se mide, y lo que se observa en el estudiante queda diferido al segundo ciclo, cuya muestra será dirigida —un curso completo, seleccionado por accesibilidad y no por procedimiento probabilístico \autocite[p.~215]{sampieri2018metodologia}— según el diseño de las recomendaciones. En su lugar, este ciclo define dos universos que recorre de forma exhaustiva. El primero es el \emph{universo de contenido}: las once unidades del canal de cálculo que el \autoref{cap:marco_teorico} desarrolla —elasticidad plana y matriz constitutiva; funciones de forma y mapeo isoparamétrico; Jacobiano y distorsión; matriz de deformación-desplazamiento; formulación débil, rigidez elemental y cuadratura de Gauss; fuerzas nodales equivalentes; ensamblaje y restricciones; resolución del sistema y recuperación de tensiones; calidad de malla; fenómenos numéricos, esto es, bloqueo, modos espurios y condicionamiento; y verificación por equilibrio y contraste con referencias— cruzadas con los cinco niveles cognitivos de la taxonomía revisada de Bloom que la tabla de especificaciones emplea \autocite{anderson2001taxonomy}. El segundo es el \emph{universo de principios}: los diez principios de aprendizaje de la \autoref{sec:fundamentos-pedagogicos}, que la matriz de trazabilidad recorre uno por uno.

Para el experimento numérico, la población sería el universo de problemas de elasticidad lineal plana —en tensión o deformación plana— resolubles con elementos cuadriláteros isoparamétricos. El conjunto de casos que sigue no es una muestra estadística de ese universo, sino una selección \emph{intencional} por valor probatorio —una muestra dirigida, no probabilística \autocite[p.~215]{sampieri2018metodologia}—, coherente con la práctica de la verificación de códigos de cálculo, donde los casos de prueba se eligen por su capacidad de ejercitar todos los términos del modelo matemático y por cubrir la topología de malla más general que el código deberá resolver, antes que por su realismo físico \autocite{oberkampf2010vv}. Cada caso aporta una referencia de naturaleza distinta y ejercita términos distintos del modelo. El método de soluciones manufacturadas provee una solución exacta arbitraria y ejercita la fuerza de volumen, las restricciones de Dirichlet no homogéneas, las dos matrices constitutivas y el mapeo sobre elementos distorsionados, en cuatro configuraciones (cuadrado unitario y cuadrilátero distorsionado, en tensión y en deformación plana). La viga de Timoshenko simplemente apoyada ejercita la carga superficial uniforme y los apoyos, y aporta las dos referencias externas: una solución analítica de la teoría de la elasticidad y un modelo de cáscara en SAP2000. La membrana de Cook ejercita la distorsión trapezoidal bajo cortante y flexión, el modo de deformación ante el que el Q4 bloquea. Entre los tres quedan cubiertos los términos del modelo de la \autoref{sec:formulacion-debil} —cargas nodales, volumétricas y superficiales uniformes, restricciones homogéneas y no homogéneas, ambos estados planos y ambos elementos—, con una excepción que se declara: la carga superficial linealmente variable no interviene en ningún caso de referencia y no tiene contraste externo; las recomendaciones proponen un caso que la incorpore. El refinamiento de malla y la comparación Q4 frente a Q9 se realizan en los dos casos con secuencia de mallas —el MMS y la membrana de Cook, ambos con $N\in\{2,4,8,16,32\}$—; la viga de Timoshenko se resuelve con una única malla fina de elementos Q9 y constituye un contraste puntual de exactitud frente a las dos referencias externas, no un estudio de convergencia.

\subsection{Procedimiento, instrumentos y reproducibilidad}
\label{sec:validacion-datos}

Cada caso numérico sigue el mismo procedimiento. Un guion construye la malla y define material, cargas y apoyos; un validador de salud del modelo —una función pura sin dependencias de la interfaz— verifica la consistencia de los datos de entrada antes de resolver: existencia de elementos, restricciones suficientes para suprimir los movimientos de cuerpo rígido, ausencia de nodos referenciados inexistentes y no degeneración de los elementos (determinante Jacobiano positivo); los errores críticos bloquean el cálculo. En la verificación por soluciones manufacturadas la consistencia de los datos está garantizada por construcción: el término fuente se deduce analíticamente de la solución exacta elegida a priori, de modo que las cargas y las restricciones de Dirichlet corresponden exactamente a ella. El sistema se ensambla y resuelve con las mismas funciones del motor que consumen los módulos educativos y la memoria de cálculo, lo que garantiza coherencia entre lo que el alumno ve explicado y lo que el programa resuelve, y las tensiones se recuperan por la secuencia descrita en el \autoref{cap:marco_teorico}.

Los instrumentos de medición numérica son los guiones de verificación y validación del proyecto (\texttt{tests/vv\_mms.py}, \texttt{tests/vv\_timoshenko.py} y \texttt{tests/vv\_cook.py}), que computan las normas de error —integradas con un punto de Gauss más por dirección que la rigidez (\autoref{sec:verificacion-validacion}) para no subestimar artificialmente el error—, los errores relativos frente a las referencias y las tasas de convergencia, y depositan los datos crudos en archivos CSV. Dentro de cada caso solo se manipulan el tipo de elemento y la densidad de malla —fijos en la viga de Timoshenko, por lo dicho en la \autoref{sec:poblacion}—; el tipo de análisis, el material, la geometría, las cargas y los apoyos, el orden de cuadratura ($2\times2$ para Q4 y $3\times3$ para Q9) y las tolerancias numéricas del motor —el cero numérico y el determinante Jacobiano mínimo admisible, constantes del programa— permanecen fijos.

Los instrumentos didácticos son cuatro y se construyen sobre la versión del software que acompaña a este documento, con las funciones que la \autoref{sec:diseno-edufem} describe y las capturas que el \autoref{anexo:manual} reproduce. El primero es la \emph{tabla de especificaciones}, que cruza las once unidades de contenido de la \autoref{sec:poblacion} con los niveles cognitivos recordar, comprender, aplicar, analizar y evaluar \autocite{anderson2001taxonomy} % LOCALIZADOR PENDIENTE: definicion de los niveles de la taxonomia revisada
y anota en cada celda el instrumento del software que la cubre —módulo educativo, memoria de cálculo, post-proceso, validador de salud, o conversión de tipo de elemento y refinamiento de la malla—; una celda se marca solo si el instrumento existe en el software distribuido y el \autoref{cap:diseno} lo describe, y queda vacía en caso contrario. Es la vía documental de la validez de contenido, la que establece que un instrumento cubre el dominio de lo que pretende medir \autocite{sampieri2018metodologia}. % LOCALIZADOR PENDIENTE: definicion de validez de contenido
El segundo es el \emph{mapa de conjeturas} \autocite{sandoval2014conjecture}, que enuncia la conjetura de alto nivel, su encarnación en herramientas y materiales, estructuras de tarea, estructuras de participación y prácticas discursivas, los procesos mediadores que esa encarnación debería producir y los resultados esperados; distingue las conjeturas de diseño, que ligan la encarnación con los procesos y son las que este ciclo contrasta, de las conjeturas teóricas, que ligan los procesos con los resultados y corresponden al segundo ciclo. El tercero es la \emph{matriz de trazabilidad}, que desarrolla la columna de encarnación del mapa: para cada principio de aprendizaje de la \autoref{sec:fundamentos-pedagogicos} registra la fuente, la decisión de diseño de la \autoref{sec:diseno-edufem} que lo encarna, la evidencia verificable en el software —una función descrita en ese capítulo, una captura del manual o una prueba de la batería de regresión— y un veredicto en tres grados: cumple, cumple parcialmente o no encarnado. El cuarto es la \emph{evaluación heurística}, que aplica las nueve dimensiones del instrumento LORI \autocite{leacock2007framework} —calidad del contenido, alineación con los objetivos de aprendizaje, retroalimentación y adaptación, motivación, diseño de la presentación, usabilidad de la interacción, accesibilidad, reusabilidad y cumplimiento de estándares— registrando para cada una la evidencia disponible y un veredicto en cuatro grados: cumple, cumple parcialmente, no evaluable sin usuarios o no aplica. La evaluación heurística es por naturaleza una inspección predictiva sin usuarios \autocite{squires1999predicting}; aquí la aplica el autor y la documenta con evidencia por ítem, lo que la convierte en una autoevaluación declarada como tal y no en un juicio de expertos independiente. Esa es su limitación, y por eso su criterio de aceptación exige evidencia y no puntajes. A estos cuatro instrumentos se suma el análisis comparativo con los antecedentes de la \autoref{tab:comparativa}, que sitúa el artefacto frente a los recursos existentes, en la tradición de la ingeniería de software educativo que integra la evaluación pedagógica en el propio desarrollo \autocite{galvis1992ingenieria,cataldi2000metodologia}.

La reproducibilidad se apoya, en lo numérico, en el determinismo del motor: ante las mismas entradas produce las mismas salidas, y cada caso se regenera por completo al reejecutar su guion con las dependencias declaradas en el repositorio. A ello se suma una batería de pruebas de regresión que verifica la estabilidad del modelo de datos ante operaciones de edición —el ciclo de transformación $\text{Q4}\rightarrow\text{Q9}\rightarrow\text{Q4}$ sin deriva numérica, la equivalencia de resultados con identificadores de nodo no contiguos y la interoperabilidad de formatos— y que se reporta en el \autoref{cap:resultados}. En lo didáctico, la reproducibilidad descansa en que cada celda de las tablas remite a una función descrita en la \autoref{sec:diseno-edufem}, a una captura del \autoref{anexo:manual} o a una prueba de la batería de regresión, de modo que cualquier lector puede rehacer cada veredicto con el software distribuido y el manual.

\subsection{Criterios de aceptación y contraste de la hipótesis}
\label{sec:validacion-resultados}

Los resultados numéricos se contrastan con referencias de naturaleza distinta y con criterios fijados a priori en los propios guiones, que los evalúan automáticamente y terminan con error si alguno falla. Para el motor: tasas de convergencia del desplazamiento observadas dentro de $\pm0{,}5$ de las teóricas ($\mathcal{O}(h^{p+1})$ en $L^2$ y $\mathcal{O}(h^{p})$ en $H^1$) en las cuatro configuraciones del MMS; para el campo de tensiones recuperado, cuya tasa teórica no está establecida para la secuencia de extrapolación y promediado nodal que emplea el motor, una cota empírica mínima declarada como tal —tasa no inferior a $1{,}4$ en Q4 ni a $1{,}9$ en Q9—; en la viga de Timoshenko, error de la flecha central inferior al $3\,\%$ frente a la solución analítica con corrección de cortante, error de la tensión normal $\sigma_x$ inferior al $1\,\%$ frente a la solución analítica y frente a SAP2000 en los tres puntos de control, y equilibrio global —suma de reacciones verticales igual a la carga total— con residuo relativo inferior a $10^{-8}$; y, en la membrana de Cook, desplazamiento del Q9 con $N=8$ dentro del $1{,}5\,\%$ del valor de referencia y flecha del Q4 inferior a la del Q9 con la misma malla, que es la firma del bloqueo por cortante. El desempeño del solucionador se reporta sin criterio de aceptación. El \autoref{cap:resultados} reporta los valores efectivamente medidos en cada magnitud, incluidas las componentes secundarias de tensión que no forman parte del criterio.

Para el potencial didáctico, los criterios se fijan sobre los tres instrumentos. Relevancia: las once unidades de contenido tienen instrumento en los niveles comprender y aplicar como mínimo; las celdas vacías en analizar y evaluar se declaran como limitaciones, no como incumplimiento. Consistencia: cada uno de los diez principios tiene al menos una encarnación con veredicto cumple o cumple parcialmente, y ningún principio queda no encarnado; además, módulo interactivo para cada una de las siete etapas elementales y de ensamblaje del canal (M1 a M7) y desarrollo con sustitución numérica de las nueve etapas en la memoria de cálculo, las tres fases del análisis operando sobre un mismo lienzo, ida y vuelta CSV/ZIP sin pérdida de entidades ni de resultados, reimportación DXF idempotente y generación de la memoria para Q4 y Q9 sin error. Practicidad esperada: ninguna dimensión aplicable del instrumento LORI queda sin evidencia, y las dimensiones con veredicto cumple parcialmente o no evaluable sin usuarios se recogen como limitaciones y como objetivos de la prueba de campo.

La hipótesis enunciada en la Introducción se contrasta cláusula por cláusula con esos criterios: la cláusula (a), relevancia, con la tabla de especificaciones (PI-1); la cláusula (b), corrección, con los criterios numéricos del motor (PI-2); la cláusula (c), consistencia, con el mapa de conjeturas y la matriz de trazabilidad, más los indicadores de cobertura, fases y formatos de la variable independiente (PI-3); y la cláusula (d), practicidad esperada, con la evaluación heurística (PI-4). La efectividad —que la exposición produzca la comprensión— no se contrasta en este ciclo: queda como conjetura teórica del mapa, y su prueba de campo se diseña en las recomendaciones con la población de la \autoref{sec:poblacion}.
